
\documentclass[11pt,a4paper]{article}
\usepackage[margin=22mm]{geometry}
\usepackage{kotex}
\usepackage{amsmath,amssymb}
\usepackage{array,longtable,booktabs}
\usepackage{xurl}
\usepackage{xcolor}
\usepackage[hidelinks]{hyperref}
\setlength{\parindent}{0pt}
\setlength{\parskip}{6pt}
\setlength{\emergencystretch}{3em}
\renewcommand{\arraystretch}{1.3}
\allowdisplaybreaks
\newcolumntype{P}[1]{>{\raggedright\arraybackslash}p{#1}}
\newcommand{\blocktitle}[1]{\par\bigskip\noindent\textbf{\large #1}\par\smallskip}

\title{overall ideas of modeling}

\begin{document}
\blocktitle{Problem of Interest}
\textbf{회사가 선정한 잠재적인 후보 도시들 중 어떤 도시에 로봇택시를 얼마나 할당할 것인가.} \\
\\
\textbf{① 기본 newsvendor 모형, \\
② 비용을 추가한 모형 \\
③ 이익을 최대화하는 모형 \\
④ multi-penalties.} \\
아래는 모두 \textbf{도시 선택과 차량 수를 동시에 결정}하는 식.

모든 수요·처리량·비용은 \textbf{동일한 평가기간} 기준으로 맞춰야 한다. 평가기간은 연구자가 정하고, 모델은 그 기간에 적용할 도시/차량 배치를 결정해.

%%%%%%%%%%%%%%%%%%%%%%%%%%%%%%%%%%%%%%%%%%%%%%%%%%%%%
\blocktitle{1. 세 모형에 공통으로 사용하는 기호}

\begin{longtable}{@{}P{0.12\textwidth}P{0.34\textwidth}P{0.46\textwidth}@{}}
\toprule
기호 & 의미 & 값을 어떻게 정하나? \\
\midrule
\endhead
$\mathcal C$ & 잠재 후보 도시 집합 & 회사의 후보 목록. 현재 회사 현황 자료는 참고 \\
$\mathcal S$ & 수요 시나리오 집합 & 연구자가 생성 \\
$y_c$ & 도시 $c$에 진입하면 1 & \textbf{최적화가 결정} \\
$q_c$ & 도시 $c$의 배치 차량 수 & \textbf{최적화가 결정} \\
$D_c^s$ & 시나리오 $s$의 로봇택시 잠재수요 & \textbf{자료로 추정 + 부족한 부분은 가정} \\
$\kappa_c$ & 차량 한 대의 기간 내 처리 가능 trip 수 & \textbf{운영자료 또는 가정}. (처음에는 고정) \\
$Q$ & 전체 배치 가능한 차량 수 & 회사 입력. 없으면 실험값 설정 \\
$\bar q_c$ & 도시별 차량 배치 상한 & 시설·운영 한도. 별도 한도가 없으면 $Q$ \\
$a_c$ & 계획 시점에 해당 도시에서 운영 가능하면 1 & 허가·ODD·시설자료로 판정. 미확인은 별도 가정 \\
$B$ & 동시에 진입할 수 있는 도시 수의 상한 & \textbf{실제 제한이 있을 때만} 회사 입력 또는 실험값 \\
$R_c^s$ & \text{실제 처리량} & 모델 내 배치 결과에 따라 계산되는 값\\
$O_c^s$ & \text{남는 처리용량} & 모델 내 배치 결과에 따라 계산되는 값\\
$U_c^s$ & \text{미처리 잠재수요} & 모델 내 배치 결과에 따라 계산되는 값\\



\bottomrule
\end{longtable}

$D_c^s$는 불확실한 수요를 표본으로 만든 값이야. \textbf{시나리오를 만든 뒤 최적화할 때는 고정 입력}이고, 모형이 수요를 선택하는 것은 아니야.

%%%%%%%%%%%%%%%%%%%%%%%%%%%%%%%%%%%%%%%%%%%%%%%%%%%%%
\blocktitle{2. 공통 제약식}

\[
\boxed{
\begin{aligned}
y_c\le q_c&\le \bar q_c y_c &&\forall c\in\mathcal C,\\
\sum_{c\in\mathcal C}q_c&\le Q, \quad \sum_{c\in\mathcal C}y_c\le B\\
y_c&\le a_c &&\forall c\in\mathcal C,\\
y_c&\in\{0,1\},\quad &&q_c\in\mathbb Z_{\ge0}.
\end{aligned}}
\]

\begin{longtable}{@{}P{0.34\textwidth}P{0.60\textwidth}@{}}
\toprule
제약 & 의미 \\
\midrule
\endhead
$y_c\le q_c\le\bar q_cy_c$ & 진입하지 않으면 0대, 진입하면 최소 1대부터 배치 \\
$\sum_cq_c\le Q$ & 전체 차량 한도 \\
$\sum_{c\in\mathcal C}y_c\le B$ & 선택할 도시 수 제한\\
$y_c\le a_c$ & 운영 가능한 도시에만 진입 \\

\bottomrule
\end{longtable}

%%%%%%%%%%%%%%%%%%%%%%%%%%%%%%%%%%%%%%%%%%%%%%%%%%%%%
\blocktitle{3. 배치 결과에 따라 계산되는 값}
\textbf{수요 $D_c^s$와 공급능력 $\kappa_cq_c$가 결정되면 배치 결과에 따라 계산되는 값}은 다음 세 가지로, 별도로 조사할 변수가 아니다.

\[
\begin{aligned}
R_c^s&=\min\{D_c^s,\kappa_cq_c\}
&&\text{실제 처리량},\\
O_c^s&=(\kappa_cq_c-D_c^s)^+
&&\text{남는 처리용량},\\
U_c^s&=(D_c^s-\kappa_cq_c)^+
&&\text{미처리 잠재수요}.
\end{aligned}
\]

여기서 $(x)^+=\max\{x,0\}$. 세 값의 단위는 모두 \textbf{trip 수}.


%%%%%%%%%%%%%%%%%%%%%%%%%%%%%%%%%%%%%%%%%%%%%%%%%%%%%%%%%%%%%%%%%
\blocktitle{① 기본형: 차량이 남는 문제와 부족한 문제를 최소화 — 먼저 만들 모형}

\[
\boxed{
\min_{y,q}\;
\frac{1}{|\mathcal S|}
\sum_{s\in\mathcal S}\sum_{c\in\mathcal C}
\left[
\underbrace{O_c^s}_{\text{남는 용량}}
+
\underbrace{\rho U_c^s}_{\text{부족에 대한 손실}}
\right]
}
\]

\begin{longtable}{@{}P{0.13\textwidth}P{0.38\textwidth}P{0.41\textwidth}@{}}
\toprule
추가 기호 & 의미 & 설정 방법 \\
\midrule
\endhead
$\rho>0$ & 부족 1건을 남는 용량 1건보다 얼마나 중요하게 보는지 & \textbf{연구자 가정 또는 회사 선호}. 예: $1,2,5$를 비교하는 민감도 실험 \\
\bottomrule
\end{longtable}

\textbf{장점:} \textcolor{red}{정확한 회사 비용자료 없이도 시작할 수 있어.}

\textbf{범위:} \textcolor{red}{진입비용을 제외한 배치 기준모형}


%%%%%%%%%%%%%%%%%%%%%%%%%%%%%%%%%%%%%%%%%%%%%%%%%%%%%%%%%%%%%%%%%%%%%
\blocktitle{② 비용형: 진입·차량 비용까지 포함해 총비용 최소화}

\[
\boxed{
\begin{aligned}
\min_{y,q}\; &
\underbrace{\sum_c F_cy_c}_{\text{도시 진입비용}}
+
\underbrace{\sum_c K_cq_c}_{\text{차량 배치·보유비용}}\\
&+
\frac{1}{|\mathcal S|}
\sum_s\sum_c
\left[
\underbrace{c_c^+O_c^s}_{\text{잔여용량 손실}}
+
\underbrace{c_c^-U_c^s}_{\text{미처리 기회 손실}}
\right]
\end{aligned}
}
\]

\begin{longtable}{@{}P{0.12\textwidth}P{0.32\textwidth}P{0.48\textwidth}@{}}
\toprule
추가 기호 & 의미·단위 & 설정 방법 / 필요한 자료 \\
\midrule
\endhead
$F_c$ & 해당 도시 진입 고정비용, 원 또는 달러 & \textbf{추가자료 필요:} 거점 구축·초기 준비 비용. 없으면 명시적 가정 \\
$K_c$ & 차량 1대의 평가기간 배치·보유비용 & \textbf{추가자료 필요:} 차량·보험·고정 운영비 등. 기존 보유차라면 과거 구입비를 무조건 다시 넣지 않음 \\
$c_c^+$ & 남는 처리용량 1 trip당 추가 손실 & 회사 자료 또는 가정. \textbf{$K_c$와 같은 비용을 중복 계산하지 않음} \\
$c_c^-$ & 놓친 잠재수요 1 trip당 손실 & 기여이익·고객 이탈 등의 자료 또는 가정 \\
\bottomrule
\end{longtable}

%%%%%%%%%%%%%%%%%%%%%%%%%%%%%%%%%%%%%%%%%%%%%%%%%%%%%%%%%%%%%%%%
\blocktitle{③ 이익형: 처리한 이용에서 얻는 이익을 최대화}

\[
\boxed{
\begin{aligned}
\max_{y,q}\; &
\frac{1}{|\mathcal S|}
\sum_s\sum_c
\underbrace{m_cR_c^s}_{\text{처리한 이용의 기여이익}}\\
&-
\underbrace{\sum_cF_cy_c}_{\text{진입비용}}
-
\underbrace{\sum_cK_cq_c}_{\text{차량 고정비용}}
\end{aligned}
}
\]

\begin{longtable}{@{}P{0.12\textwidth}P{0.38\textwidth}P{0.42\textwidth}@{}}
\toprule
추가 기호 & 의미 & 설정 방법 / 필요한 자료 \\
\midrule
\endhead
$m_c$ & trip 1건의 수입에서 그 trip에 따라 발생하는 변동비를 뺀 금액 & \textbf{추가자료 필요:} 운임, 운행·에너지 등 변동비. 없으면 가정 \\
\bottomrule
\end{longtable}

여기서는 고객을 처리하지 못하면 받을 이익이 줄어들어. \textbf{같은 매출 손실을 $c_c^-U_c^s$로 다시 더하면 중복될 수 있어.}

회사 관점은 가장 직접적이지만, \textcolor{red}{\textbf{운임·변동비 자료가 필요}}해. 수익성이 없으면 어느 도시에도 진입하지 않는 결과도 가능해.



\end{document}
